Being able to track position has become increasingly important, and is something that is often taken for granted.
We are used to being able to tell where we are at all times and the most common way to do this is through the use of GPS.
However GPS are not always available, and there are many situations where GPS is not a viable option.
For example if you are inside a building or underground, GPS signals may not be able to reach you.
In these situations it's of interest to look at alternative methods for tracking position.

GPS, global positioning system, is a satellite-based navigation system that allows users to determine their exact location anywhere on Earth.
GPS was originally developed by the United States Department of Defense for military purposes, but has since been made available for civilian use \cite{GPS}.
The system consists of a network of satellites that orbit the Earth and transmit signals to GPS receivers on the ground.
Based on the time it takes for the signals to reach the receiver, the receiver can calculate its distance from each satellite and use this information to determine its location.
The system is highly accurate and can provide location information with an accuracy of a few meters.
The GPS system is however not fit for determining the position inside of a building as the signals from the satellites are often blocked by walls and other obstacles.


